\subsection{代数扩张的传递性。在扩域中相对代数闭的域}

命题3. --- 设E和F是域K的两个扩域，使得 \(K \subset E \subset F\) 。要使 F 在 K 上是代数的，当且仅当 E 在 K 上是代数的且 F 在 E 上是代数的。

该条件是必要的，由 V，第 17 页，推论 2 可知。下面证明其充分性。设 x 为 F 的任意元素；它在 E 上是代数的；设 \(g \in E[X]\) 为其在 E 上的极小多项式。若 A 为 g 的系数的（有限）集合，则 \(g \in K(A)[X]\) ，因此 x 在 \(K(A)\) 上是代数的，且 \(K(A \cup \{x\}) = K(A)(x)\) 在 \(K(A)\) 上是有限次的。现在 \(A \subset E\) 且E在K上是代数的，因此 \(K(A)\) 在K上是有限次的（由定理2）。因此（第V卷，第10页，定理1）， \(K(A \cup \{x\})\) 在K上是有限次的，这表明x在K上是代数的（V，第18页，命题2）。

定义2. —— 域E的子域K称为在E中相对代数闭的，如果E中每个在K上代数的元素都属于K。

这等价于说，K 是包含在 E 中的唯一的 K 的代数扩张。每个域K在自身中都是相对代数闭的。在第4节中，我们将研究在每个扩域中都是相对代数闭的域。

命题4. --- 设E是域K的一个扩域；E中在K上代数的元素之集L构成E的一个子扩域，且L在E中相对代数闭。

因为域 \(\mathbf{K}(\mathbf{L})\) 在 \(\mathbf{K}\) 上是代数的（定理2的推论1），因此 \(\mathbf{K}(\mathbf{L}) \subset \mathbf{L}\) ；由此可得 \(\mathbf{K}(\mathbf{L}) = \mathbf{L}\) ，且 \(\mathbf{L}\) 是一个域。另一方面，如果 \(x \in \mathbf{E}\) 在 \(\mathbf{L}\) 上是代数的，则它在 \(\mathbf{K}\) 上也是代数的（命题3），因此属于 \(\mathbf{L}\) 。

K在E中的相对代数闭包是指E中所有在K上代数的元素组成的K的扩张L；它是包含在E中的K的最大代数扩张。

